\begin{problem}{Zouari's Modular Strategy (Easy Version)}{standard input}{standard output}{1 second}{256 megabytes}

\begin{itemize}
  \item \textbf{Easy version}: You only need to check if it is possible to transform the \textbf{unordered multiset} $A$ such that $(A \bmod k) = B$. If it is possible, print \textbf{YES}; otherwise, print \textbf{NO}.
  \item \textbf{Hard version}: In addition to deciding whether it is possible, you must also output the minimal number of operations required to achieve it. 

\textbf{NB : } The constraints of the two versions are different.
\end{itemize}


Mohamed Zouari, the brilliant engineer from Sfax, Tunisia, dedicated his life to teaching Palestinian students about \textbf{advanced drone technology} and strategic problem-solving. During one of his lessons, he wanted to emphasize planning, patience, and incremental progress.

On the board, he wrote a challenge for the students. The challenge has two versions:


You are given three integers $n$, $k$, and $y$, and two \textbf{unordered multisets}
$$A = \{a_1, a_2, \dots, a_n\}\quad\text{and}\quad B = \{b_1, b_2, \dots, b_n\}.$$

You are allowed to perform the following operation any number of times (possibly zero):
\begin{itemize}
  \item Select an index $i$ ($1 \le i \le n$) and perform $a_i := a_i + y$.
\end{itemize}

\textbf{Your goal} is to determine whether you can  make $(A \bmod k) = B$ or not.

We say that $A \bmod k = B$ if and only if the following condition holds:
$$
\{\, a_i \bmod k \mid 1 \le i \le n \,\}
=
\{\, b_i \mid 1 \le i \le n \,\},
$$
where both sides are considered as \textbf{unordered multisets}. In other words, after taking each element of $A$ modulo $k$, the resulting multiset of remainders is exactly the same as $B$, ignoring the order of elements but taking into account their multiplicities.For example, the resulting array ($A$ mod $k$) $= [2,3,1,2,4]$ and the resulting $B = [2,2,3,4,1]$ are considered equal, but the resulting array ($A$ mod $k$) $= [1,2,2]$ and $ B =[1,2,3]$  are not equal ($2$ appears twice in the array $A$ while $3$ is absent).

\InputFile
The first line contains a single integer $t$ $(1 \le t \le 10^4)$ --- the number of test cases.

The first line of each test case contains three integers $n$, $k$, and $y$ ($1 \le n \le 2 \cdot 10^5$, $1 \le k \le 10^9$, $0 \le y \le 10^9$).

The second line of each test case contains $n$ integers $a_1, a_2, \dots, a_n$ ($0 \le a_i \le 10^9$).

The third line of each test case contains $n$ integers $b_1, b_2, \dots, b_n$ ($0 \le b_i \le 10^9$).

It is guaranteed that the sum of $n$ over all test cases does not exceed $2 \cdot 10^5$ .

\OutputFile
Print \textbf{YES} if it is possible to obtain $(A \bmod k) = B$ using the allowed operation; otherwise, print \textbf{NO}.

\Example

\begin{example}
\exmpfile{example.01}{example.01.a}%
\end{example}

\end{problem}

